\begin{ledgroupsized}[r]{120mm}
\footnotesize
\pstart
\noindent\textbf{Überlieferung:}
\pend
\end{ledgroupsized}
\begin{ledgroupsized}[r]{114mm}
\footnotesize
\pstart
\parindent-6mm
\makebox[6mm][l]{\textit{L}}%
Konzept:
LH~XXXV~10, 16~Bl.~3\textendash4. 
Ein Bogen~2\textsuperscript{o};
\href{\Scan{LH_35_10_16_003-004wz}}%
{Wasserzeichen} im Falz.
Eine Seite auf Bl.~4~v\textsuperscript{o};
Bl.~3 bis 4~r\textsuperscript{o} überliefern die Auszüge N.\;\ref{60123}.
Der Text und die Zeichnungen sind gestrichen.
\pend
\end{ledgroupsized}
\selectlanguage{latin}
\vspace{8mm}
\footnotesize
\pstart\noindent
\href{\Scan{LH_35_10_16_004v}}%
{\lbrack4~v\textsuperscript{o}\rbrack}
\edtext{Inspectione}{\lemma{}\Bfootnote{Inspectione \textit{erg.~L}}} 
\edtext{\cite{00541}\textit{Neostaticae}\indexS{Neostatica} celeberrimi viri Hieronymi Saccherii\indexN{\textls{Saccheri}, Girolamo 1667\textendash 1733}}{%
\lemma{\textit{Neostaticae} celeberrimi viri Hieronymi Saccherii}%
\Cfootnote{\indexN{\textls{Saccheri}, Girolamo 1667\textendash 1733}\textsc{G. Saccheri}, \cite{00541}\textit{Neo-statica}, Mailand 1708.}}
mirifice sum delectatus. Animadverti enim habere virum meditationes profundas, 
atque 
\edtext{ipsius etiam calculi infinitesimalis\indexS{calculus infinitesimalis} 
\edtext{a me nonnihil promoti}{\lemma{}\Bfootnote{a me nonnihil promoti \textit{erg.~L}}} 
veras rationes inspexisse, ut apparet 
ex p.~79.}{\lemma{ipsius \lbrack...\rbrack\ p.~79}\Cfootnote{%
\cite{00541}a.\,a.\,O., lib.~3, Praenotanda, Nr.~5, S.\;79. Dort verweist \indexN{\textls{Saccheri}, Girolamo 1667\textendash 1733}Saccheri außerdem auf einige \glqq clarissimos Geometras\grqq, allen voran \glqq Liebnitium\grqq\ (\textit{sic!}).}}
Et vellem 
\edtext{omnia quae dedit habere et eo \lbrack cui\rbrack\ par}{\lemma{omnia}\Bfootnote{%
\textit{(1)}~eo quo %
\textit{(2)}~quae dedit
\textit{(a)}~et
\textit{(b)}~habere et eo \bars{quo \textit{ändert Hrsg.}} par~\textit{L}}} 
est studio expendere \edlabel{LH_35_10_16_004_1a}liceret.
\edtext{}{{\xxref{LH_35_10_16_004_1a}{LH_35_10_16_004_1b}}\lemma{liceret.}\Bfootnote{%
\textit{(1)}~Pulchre \bars{etiam \textit{erg.}} uti videtur tangente cujusdam curvae ad circuli tetragonismum. %
\textit{(2)}~Haud dubie \lbrack...\rbrack\ principiis. \bars{Pulchre \lbrack...\rbrack\ circuli. \textit{erg.}} 
Nescio an viderit \textit{propositionem} %
\textit{(a)}~meam %
\textit{(b)}~\textit{generalem compositionis motuum} %
\textit{(aa)}~.~Esto mobile %
\textit{(bb)}~ante \lbrack...\rbrack\  productam. %
\textit{(aaa)}~Pon %
\textit{(bbb)}~Mobile \textit{M} %
\textit{(aaaa)}~habeat eodem %
\textit{(bbbb)}~feratur motu composito ex motibus
\textit{(aaaaa)}~ponentibus
\textit{(bbbbb)}~componentibus quotcunque \bars{\textit{MF}, \textit{MG}, \textit{MH} \textit{erg.}}. Fingatur ferri unoquoque %
\textit{(aaaaa-a)}~et intra ta
\textit{(bbbbb-b)}~via centr %
\textit{(ccccc-c)}~separatim, tanquam essent totidem mobilia 
\textit{(aaaaa-aa)}~seu
\textit{(bbbbb-bb)}~seu \lbrack...\rbrack\ commune %
\textit{(aaaaa-aaa)}~trium %
\textit{(bbbbb-bbb)}~horum \lbrack...\rbrack\ feretur \bars{mobile \textit{M} \bars{recta \textit{MK} \textit{gestr.}} \textit{erg.}} via hujus centri, %
\textit{(aaaaa-aaaa)}~et %
\textit{(bbbbb-bbbb)}~\bars{seu directione \textit{erg.}} \textit{MK}, et velocitate 
\bars{ut \textit{ML} \textit{erg.}} quae sit ad 
\bars{\textit{MK} \textit{erg.}} progressum centri gravitatis 
\bars{locorum \textit{erg.}} communis, \lbrack...\rbrack\ 3 ad 1. %
\textit{(aaaaa-aaaaa)}~In uno tamen cursoria mea lectio adhaesit. %
\textit{(bbbbb-bbbbb)}~\lbrack/\rbrack~In uno tamen \bars{adhaesit \textit{erg.}} cursoria~\textit{L}}}%
Haud dubie enim non parum inde proficerem. Sane ingeniose admodum et profunde versari videtur 
in ipsis rei phoranomicae\indexS{res phoranomica} demonstrandis principiis.\indexS{principia rei phoranomicae}
Pulchre etiam videtur uti tangente cujusdam curvae\indexS{tangens cujusdam curvae} ad tetragonismum circuli.\indexS{tetragonismus circuli} 
\edlabel{35_10_16_004_4a}%
Nescio an viderit 
\edtext{\textit{propositionem generalem compositionis motuum}\indexS{propositio generalis compositionis motuum}\indexS{compositio motuum} ante multos annos a me in \textit{diario eruditorum} Gallico productam.}{\lemma{propositionem \lbrack...\rbrack\ productam}%
\applabel{FnC:35_10_16_004_regle}%
\Cfootnote{\textsc{G.\,W.~Leibniz}, \cite{00542}\textit{Regle generale de la composition des mouvemens}, \textit{JS} vom 7.~September 1693, S.\;417\textendash419.}}
Mobile \textit{M} feratur motu composito ex motibus componentibus\indexS{motus componens} quotcunque\indexS{motus ex motibus componentibus quotcunque compositus} \textit{MF}, \textit{MG}, \textit{MH}. 
Fingatur ferri unoquoque separatim,\indexS{mobile latum unoquoque motu separatim} tanquam essent totidem mobilia seu tanquam simul veniret in \textit{F}, \textit{G}, \textit{H}. 
Quaeratur centrum commune\indexS{centrum commune} horum trium punctorum seu locorum quae unoquoque motu componente\indexS{motus componens} intra idem tempus nancisceretur, 
feretur mobile\indexS{mobile} \textit{M} via hujus centri,\indexS{via centri communis} seu directione \textit{MK}, et velocitate ut \textit{ML} 
quae sit ad \textit{MK} progressum centri gravitatis\indexS{progressus centri gravitatis} locorum communis,\indexS{centrum gravitatis commune} 
ut numerus motuum componentium\indexS{numerus motuum componentium} est ad unitatem, seu hic ut 3 ad 1.%
\edlabel{35_10_16_004_4b}%
\pend 
\begin{edfigure}%
\normalsize
\hspace*{2em}%
\includegraphics[scale=1]{%
gesamttex/edit_VIII-4/images/LH_35_10_16_004_d1_004v.pdf}
\edcaption{\edtext{\lbrack\textit{Fig.~1, gestr.}\rbrack}{%
\lemma{\lbrack\textit{Fig.~1}\rbrack}%
\Cfootnote[nonum]{Zwei ebenfalls gestrichene Entwürfe zum Diagramm werden nicht wiedergegeben.}}}
\end{edfigure}
\pstart
\edlabel{35_10_16_004_5a}%
In uno tamen adhaesit cursoria\edlabel{LH_35_10_16_004_1b}
mea lectio\lbrack,\rbrack\ nam 
\edtext{satis attenta}{\lemma{satis}\Bfootnote{%
\textit{(1)}~dilige %
\textit{(2)}~attenta~\textit{L}}} 
esse non potuit: video asseri 
\edlabel{LH_35_10_16_004_2a}%
\edtext{}{%
\xxref%
{LH_35_10_16_004_2a}{LH_35_10_16_004_2b}%
\lemma{lib.~2. prop.~15. corollar.}\Cfootnote{\indexN{\textls{Saccheri}, Girolamo 1667\textendash 1733}\textsc{Saccheri}, \cite{00541}\textit{Neo-statica}, S.\;69.}}%
lib.~2. prop.~15.
\edtext{corollar.\edlabel{LH_35_10_16_004_2b} solicitationes gravium\indexS{solicitatio gravis} ex quiete}{\lemma{corollar.}\Bfootnote{%
\textit{(1)}~impetus %
\textit{(2)}~conatus gravium ex quie %
\textit{(3)}~solicitationes gravium ex quiete~\textit{L}}} 
versus centrum\indexS{solicitationes gravium versus centrum} gravium commune\indexS{centrum gravium commune} 
non posse esse in quavis distantia aequales,\indexS{solicitationes in quavis distantia aequales} et multo minus in majore 
\edtext{distantia posse, sed potius in majore distantia debere esse}{\lemma{distantia}\Bfootnote{%
\textit{(1)}~\bars{posse \textit{erg.}} esse minores, %
\textit{(2)}~posse, sed \bars{potius \textit{erg.}} in majore distantia \bars{debere \textit{erg.}} esse~\textit{L}}} 
majores; mihi tamen contrarium persuasum fuit. 
\edtext{Quin olim in \textit{tentamine Motuum coelestium}	%
posui 
\edtext{esse solicitationes paracentricas\indexS{solicitatio paracentrica} reciproce}{\lemma{esse}\Bfootnote{%
\textit{(1)}~attractiones %
\textit{(2)}~solicitationes paracentricas
\textit{(a)}~in
\textit{(b)}~reciproce~\textit{L}}} 
ut 
\edtext{quadrata distantiarum;\indexS{quadratum distantiae}\indexS{solicitationes paracentricae reciprocae quadratis distantiarum}}{\lemma{quadrata}\Bfootnote{%
\textit{(1)}~distantiae %
\textit{(2)}~distantiarum;~\textit{L}}} 
eo argumento quo statuunt dudum Optici\indexS{Optici}\lbrack,\rbrack\ illuminationes\indexS{illuminatio}\indexS{illuminationes reciprocae quadratis distantiarum ab illuminante} esse reciproce ut quadrata distantiarum\indexS{quadratum distantiae} ab \indexS{illuminans}illuminante.%
\edlabel{35_10_16_004_5b}%
}{\lemma{Quin \lbrack...\rbrack\ illuminante}%
\applabel{FnC:35_10_16_004_Tentamen}%
\Cfootnote{Leibniz spielt auf sein \cite{00543}\textit{Tentamen de motuum coelestium causis} an, das in den \textit{Acta Eruditorum} von Februar 1689 erschien (S.\;82\textendash96). 
Die einschlägige Stelle findet sich allerdings nicht in der veröffentlichten Fassung, sondern nur in der späteren sog. \glqq zweiten Bearbeitung\grqq: 
siehe bspw. 
LH~XXXV~9, 2~Bl.~54~v\textsuperscript{o} und \cite{01043}\textit{LMG}~VI, S.\;165\,f.}} 
\pend
\begin{edfigure}%
\normalsize
\includegraphics[scale=1]{%
gesamttex/edit_VIII-4/images/LH_35_10_16_004_d2_004v.pdf}
\edcaption{%
\edtext{\lbrack\textit{Fig.~2, gestr.}\rbrack}{%
\lemma{\lbrack\textit{Fig.~2}\rbrack}%
\Cfootnote[nonum]{Vorlage ist die Figur aus \indexN{\textls{Saccheri}, Girolamo 1667\textendash 1733}\textsc{Saccheri}, \cite{00541}\textit{Neo-statica}, S.\;67.}}}
\end{edfigure}
\pstart
Rem autem 
\edtext{possibilem esse, ut solicitationes\indexS{solicitatio} decrescant aucta distantia,\indexS{solicitatio aucta distantia decrescens} sic ostendo.}{\lemma{possibilem}\Bfootnote{%
\textit{(1)}~, saltem sic ostendo. %
\textit{(2)}~esse ut solicitationes decrescant aucta distantia, sic ostendo.~\textit{L}}} 
\edtext{Esto libra\indexS{libra} \textit{CNR}, centrum aequilibrii\indexS{centrum aequilibrii} simul}{\lemma{Esto}\applabel{FnB:LH_35_10_16_004_esto-tecum}\Bfootnote{%
\textit{(1)}~(tecum) %
\textit{(2)}~ponamus pon %
\textit{(3)}~libra \textit{CNR}, %
\textit{(a)}~punctum aequilibrii si %
\textit{(b)}~centrum aequilibrii simul~\textit{L}}} 
et motus\indexS{centrum motus} 
\edtext{sit}{\lemma{}\Bfootnote{sit \textit{erg.~L}}} 
\textit{N}. Ponamus pondera 
\edtext{in \textit{C} et \textit{R}}{\lemma{in}\Bfootnote{%
\textit{(1)}~\textit{R} et %
\textit{(2)}~\textit{C} et \textit{R}~\textit{L}}} 
esse aequalia sed eorum solicitationes\indexS{solicitatio}\indexS{solicitationes gravium versus centrum} versus centrum commune gravium\indexS{centrum commune gravium} \textit{D} esse inaequales. 
\edtext{Pro manifesto semper habui}{%
\lemma{Pro manifesto semper habui}%
\Cfootnote{Siehe bspw. die \textit{Solutio problematis R.\,P.\;Vannii S.\,J.} (April bis Juni 1690), N.\;\ref{41823} in diesem Band, bes. \refpassage{35_10_13_010-011_17a}{35_10_13_010-011_17b}.}}
actis ex \textit{N} in \textit{CD}, \textit{RD}, normalibus \textit{NL}, \textit{NX} fore aequilibrium\indexS{aequilibrium} si sint 
\edtext{solicitationes\indexS{solicitatio} reciproce}{%
\lemma{solicitationes}%
\Bfootnote{%
\textit{(1)}~in e
\textit{(2)}~reciproce~\textit{L}}}
ut perpendiculares,\indexS{solicitationes reciprocae perpendicularibus} seu \textit{NL} ad \textit{NX}, ut solicitatio\indexS{solicitatio} in \textit{R} ad solicitationem 
\edtext{in \textit{C}. Nam fingi}{\lemma{in \textit{C}.}\Bfootnote{%
\textit{(1)}~Idque comperi %
\textit{(2)}~Nam fingi~\textit{L}}} 
potest \textit{NL}, \textit{NX}, esse lineas rigidas, ubi
patet idem esse, trahas in \textit{L} an in \textit{C}. Video etiam hoc 
\edlabel{LH_35_10_16_004_3a}%
\edtext{}{%
\xxref%
{LH_35_10_16_004_3a}{LH_35_10_16_004_3b}%
\lemma{consentire}\applabel{FnB:LH_35_10_16_004_consentire-tuae}\Bfootnote{%
\textit{(1)}~tuae %
\textit{(2)}~propositioni~\textit{L}}}%
consentire
\edtext{propositioni\edlabel{LH_35_10_16_004_3b}
sextae libri secundi,}{\lemma{propositioni sextae libri secundi}\Cfootnote{\indexN{\textls{Saccheri}, Girolamo 1667\textendash 1733}\textsc{Saccheri}, \cite{00541}\textit{Neo-statica}, S.\;59\,f.}}
nam \textit{NL} est ad \textit{NX} in ratione composita \textit{CN} ad \textit{RN} et \textit{RD} ad \textit{CD}. Hinc si 
\edtext{jam}{\lemma{}\Bfootnote{jam \textit{erg.~L}}} 
\edtext{ponamus cum cl\textsuperscript{mo} Autore}{\lemma{ponamus}\applabel{FnB:LH_35_10_16_004_ponamus-tecum}\Bfootnote{%
\textit{(1)}~tecum %
\textit{(2)}~cum cl\textsuperscript{mo}  Autore~\textit{L}}} 
dari puncta \textit{R}, \textit{N}, \textit{D}, 
\edtext{angulo \textit{RND} recto\lbrack,\rbrack}{\lemma{}\applabel{FnB:LH_35_10_16_004_angulo-recto}\Bfootnote{angulo \textit{RND} recto \textit{erg.~L}}} 
dari etiam solicitationes\indexS{solicitatio} \textit{R} et \textit{C}, quaerique punctum \textit{C}, 
\edtext{seu rectam \textit{NC}}{%
\lemma{}%
\Bfootnote{%
seu rectam \textit{NC} 
\textit{erg.~L}}}
ut \textit{N} sit centrum motus\indexS{centrum motus} et aequilibrii.\indexS{centrum aequilibrii} Dico 
\edtext{debere in \textit{RN} quantum satis producta punctum \textit{C} assumi}{\lemma{debere}\Bfootnote{%
\textit{(1)}~\textit{C} assu %
\textit{(2)}~in \textit{RN} quantum satis producta punctum \textit{C} assumi~\textit{L}}} 
tale, ut sit \textit{NL}, ad datam \textit{NX}, in ratione solicitationis\indexS{solicitatio} \textit{R} ad solicitationem \textit{C}. 
Sed hinc patet subesse quosdam limites datis punctis \textit{R}, \textit{N}, \textit{D}. 
Nam ratio solicitationis\indexS{solicitatio} in \textit{R} ad solicitationem in \textit{C}, 
potest esse quantumlibet parva, quia puncto \textit{C} satis admoto ad \textit{N}, etiam \textit{NL} fiet utcunque 
parva, sed non potest ratio solicitationis in \textit{R}, ad solicitationem 
\edtext{in \textit{C} assumi}{\lemma{in \textit{C}}\Bfootnote{%
\textit{(1)}~esse %
\textit{(2)}~assumi~\textit{L}}} 
major quam ratio \textit{ND} ad \textit{NX} data, quia utcunque removeas \textit{C} ab \textit{N}, seu augeas \textit{NC}, nunquam tamen \textit{NL} perveniet ad ipsam \textit{ND}. Ponamus jam 
\edtext{solicitationes paracentricas,\indexS{solicitatio paracentrica}}{\lemma{}\Bfootnote{solicitationes \bars{gravium \textit{gestr.}} paracentricas,~\textit{L}}} 
esse reciproce ut distantias,\indexS{solicitationes paracentricae reciprocae distantiis} seu minui ratione distantiarum, 
fiat solicitatio in \textit{R} ad solicitationem in \textit{C}, vel breviter \textit{R} ad \textit{C}, ut \textit{CD} ad \textit{RD}, ex hypothesi. 
Rursus \textit{R} est ad \textit{C} ut \textit{NL} ad \textit{NX} ex principiis mechanicis,\indexS{principia mechanica}
ergo fit \textit{CD} ad \textit{RD}, ut \textit{NL} ad \textit{NX} sed (ex Geometria)\indexS{Geometria} 
\textit{NL} est ad \textit{NX} in ratione composita \textit{CN} ad \textit{RN} et \textit{RD} ad \textit{CD}. 
Ergo fiet \textit{CN} ad \textit{RN} in duplicata ratione \textit{CD} 
\edtext{ad \textit{RD}. Quodsi (ut ego ex vero faciendum arbitror)}{\lemma{ad \textit{RD}.}\Bfootnote{%
\textit{(1)}~Quod 
\textit{(2)}~F
\textit{(3)}~Quodsi (ut ego
\textit{(a)}~arbitro %
\textit{(b)}~ex vero faciendum %
\textit{(aa)}~arb %
\textit{(bb)}~ar %
\textit{(cc)}~arbitror)~\textit{L}}} 
ponatur \textit{R} ad 
\edtext{\textit{C}, ut quadrat. \textit{CD} ad quadrat. \textit{RD}, manifestum}{\lemma{\textit{C}, ut}\Bfootnote{%
\textit{(1)}~quadr.
\textit{(2)}~quadrat. \textit{CD} ad quadrat. \textit{RD},
\textit{(a)}~eode %
\textit{(b)}~eadem %
\textit{(aa)}~ratioci %
\textit{(bb)}~ratio %
\textit{(cc)}~sim %
\textit{(c)}~manifestum~\textit{L}}} 
est simili ratione, fieri \textit{CN} ad \textit{RN}, in triplicata ratione \textit{CD} ad \textit{RD}\lbrack,\rbrack\ in quibus omnibus nihil est impossibilitatis.
\pend 
\normalsize
